\chapter{AI-Assisted Formal Development}
\label{ch:ai-development}

The chapters up to this point have described a manufacturing pipeline
whose central move is treating every candidate producer as untrusted
with respect to the formal corpus, regardless of what that producer is.
This chapter turns the same discipline onto the process that actually
built the pipeline reported in this thesis: Claude Code and related
language-model tooling were used, throughout the run described in
Chapter~\ref{ch:procint} and Chapter~\ref{ch:empirical}, to generate
candidate patches, proof sketches, ontology fragments, documentation
drafts, and release metadata. This chapter asks what standing that
assistance has, what structural guardrails exist to keep it from
acquiring standing it has not earned, and what happened the one
documented time a gap in those guardrails let a claim through anyway.
It also surveys two unledgered hand-authored surfaces --
\texttt{pylab/} and \texttt{procint/Playground/} -- that exist
precisely so that AI-assisted and human experimentation has somewhere
to happen \emph{without} implicitly claiming release standing, and
closes with a case study in this project's own self-correction
discipline: a five-rail audit that found a real gap and the mechanized
scanner built in direct response to it.

\section{The Untrusted-Producer Doctrine, Applied to Its Own Construction}
\label{sec:ai-doctrine}

paper/main.tex's disclosure of generative-AI tool use states the
governing separation directly: language-model outputs generated during
the manufacturing run were treated throughout as untrusted candidate
artifacts, admitted into \textsf{procint} only when accepted by the Lean
kernel, the sorry census, the axiom audit, schema validation against the
\texttt{procint:} ontology, and the negative-fixture and release-gate
checks developed in Chapter~\ref{ch:mfact} and Chapter~\ref{ch:procint}.
No language-model output is part of the trusted base; that base is the
Lean kernel itself, Mathlib and CSLib~\cite{mathlib2020,cslib2026} at
their pinned revisions, and the \texttt{lake}/\texttt{ggen} toolchain.
The author of the accompanying paper reviewed and accepts responsibility
for all submitted content, code, citations, and claims in both the paper
and the repository -- a human accountability statement that this
chapter's discussion of automated tooling does not dilute: nothing
described below is offered as a substitute for that review, only as
machinery that makes the review's job more tractable and its omissions
more detectable.

This separation is not a precaution taken in the abstract. General-purpose
language models have a documented Lean~3/Lean~4 syntax-confusion failure
mode, generating deprecated Lean~3 constructs despite explicit Lean~4
instructions, with GPT-4-class and substantially smaller models reported
at comparable, low accuracy on the MiniF2F-Valid
benchmark~\cite{langconfusion2025}. Two families of approach trade some
of the resulting untrusted-candidate distance for a higher raw acceptance
rate: fine-tuning a general-purpose model toward Lean~4
specifically~\cite{theoremllama2024}, or pairing a language model with
solver-guided refinement~\cite{apollo2025}. \textsf{mfact} does not take
either trade in the run this thesis reports: every candidate, however
produced, passes through the same kernel-and-audit gate regardless of
its source's measured reliability, which is precisely why the framework
developed in Chapter~\ref{ch:mfact} needed to be built at all rather than
substituted for by a higher-precision generator. A more capable producer
changes how often a candidate is admitted; it does not change whether
admission, rather than production, is what confers standing.

\section{Genetic Tactic Search: Kernel-Accept Rate as the Only Fitness Signal}
\label{sec:ai-genetic-search}

One concrete, currently exploratory piece of tooling makes the
untrusted-producer doctrine's consequences for automated search
explicit. \texttt{docs/genetic-tactic-search.md} records a
genetic-programming loop, modeled on the TPOT2 pipeline-search paradigm
of evolving and selecting structured candidates by mutation and
crossover, applied here to sequences of Lean tactics rather than
scikit-learn pipeline stages. At the time that design note was written,
two declarations in \texttt{research/wfnet/obligations.toml} --- the
crown-jewel short-circuit soundness equivalence and the unfolding
correctness statement --- were both \texttt{status = "stated"}, and the
note is explicit that closing either requires new lemmas, not
recombination of existing tactics, so the search does not target them
directly. The crown-jewel obligation has since been discharged (its
current classification, \CrownJewelStatus, is rendered below); the
warm-up strategy the design note describes remains apt regardless, and
today applies to whatever open obligations remain, currently the single
declaration discussed in
Section~\ref{sec:mfact-boundary} of Chapter~\ref{ch:mfact}.

\begin{center}
% GENERATED FILE — DO NOT EDIT BY HAND (source: quadrature graph via ProcInt.crownJewel_status)
The release-time classification of the short-circuit soundness
characterization is \textbf{\CrownJewelStatus} (ground truth:
\texttt{ProcInt.crownJewel\_status}); this fragment is generated from the
catalog and cannot report it as proven while the obligation is open.

\end{center}

The search loop itself runs entirely outside Lean, in
\texttt{scripts/genetic\_tactic\_search.py}: each candidate genome is a
sequence of tactic calls drawn from a small, hand-vetted vocabulary ---
\texttt{aesop}, \texttt{simp}, \texttt{omega}, \texttt{exact},
\texttt{apply}, \texttt{rcases}, \texttt{constructor}, \texttt{intro},
\texttt{decide} --- spliced into a scratch \texttt{.lean} file and
checked with \texttt{lake env lean}. \textbf{Fitness is kernel
accept/reject, a hard, computed signal, never an LLM self-report or a
heuristic estimate}: nothing in this tool claims a proof exists without
the Lean kernel confirming it, following the same computed-not-asserted
invariant the release manifest itself obeys. A secondary, shaping term
in the fitness function draws on the empirical finding that tactics
conforming to expert proof patterns are more likely to advance a proof
successfully than arbitrary tactic
choices~\cite{tacticexpertalign2026}; this term only breaks ties among
sequences that already type-check, or biases mutation toward more
plausible candidates -- it never substitutes for the primary
kernel-accept signal. This is a modest but real precedent: native
LLM-tactic integration into Lean~4 has reached substantial per-step
automation rates elsewhere (Lean Copilot reports 74.2\% step-level
automation against a 40.1\% \texttt{aesop}-only
baseline~\cite{leancopilot2024}), which motivates treating expert-pattern
conformance as plausible signal, without licensing the tool to call an
LLM directly for tactic proposal in its first version -- a restraint the
design note attributes in part to the syntax-confusion risk already
discussed in Section~\ref{sec:ai-doctrine}~\cite{langconfusion2025}.
Closing genuinely new theorems, as opposed to routine proof steps,
remains hard even for purpose-built systems: HyperTree Proof Search
reports 82.6\% on Metamath with online
training~\cite{hypertree2022}, and HOList's benchmark targets a large but
bounded tactic library rather than open-ended lemma
discovery~\cite{holist2019}. Both results are cited in the design note as
the reason the tool's v1 target is warm-up lemmas in
\texttt{procint/ProcInt/Playground/TacticSearchWarmup.lean}, not either
crown-jewel-class obligation directly. Infrastructure the design note
surveys as relevant context but does not depend on in v1 includes
Lean4Lean, an independently verified Lean~4
typechecker~\cite{lean4lean2024}; Lean-auto, an interface to external
automated theorem provers~\cite{leanauto2025}; LeanExplore, a semantic
and lexical search engine over Lean declarations useful for premise
selection if the tool grows beyond a fixed vocabulary~\cite{leanexplore2025};
and large synthetic Mathlib4-derived theorem corpora for training
proof-synthesis models at scale~\cite{mathlib4dataset2025}.

The candidate-to-admission boundary for this tool is stated as
precisely as the rest of the pipeline demands. A kernel-accepted tactic
sequence the script finds is a legitimate \texttt{Mfact.Candidate} with
\texttt{producer = "genetic\_tactic\_search"}, in the sense of
Section~\ref{sec:mfact-candidates}. Promotion to \texttt{Admitted} --
flipping a TTL fragment's \texttt{status} from \texttt{"stated"} to
\texttt{"proven"} and re-running the full \texttt{just render \&\& just
build \&\& just audit \&\& just manifest \&\& just certify} chain -- is a
human-reviewed, separate step; the script itself never writes to
\texttt{packs/*/fragments/*.ttl}. Success for the search is defined
narrowly, as \texttt{lake env lean <candidate-file>} exiting~0, logged
verbatim; the design note is explicit that no claim of a
\textsc{Stated}-to-\textsc{Proven} transition is made without the full
\texttt{just check} / \texttt{just manifest \&\& just certify} chain
passing on a candidate a human has promoted. The tool's own status
line is honest about its place in the larger system: exploratory,
off-ledger tooling, not part of the certified pipeline.

\section{Unledgered Hand-Authored Surfaces: \texttt{pylab/} and \texttt{procint/Playground/}}
\label{sec:ai-unledgered}

Section~\ref{sec:mfact-ledger} of Chapter~\ref{ch:mfact} introduced the
ledger law's two permanently excluded surfaces by name;
this section develops why they exist and what governs them.
\texttt{procint/Playground/} is the pre-existing hand-authored, unledgered
Lean demonstration surface: worked, runnable instances of algorithms the
repository cares about, proving that \textsf{procint}'s own formal
content is actually exercised, not merely declared, without those
demonstrations themselves acquiring release standing. \texttt{pylab/} is
the Python-language counterpart, existing for the same reason but for
tools that have no reasonable Lean formalization: AutoML pipeline
search, the \texttt{pm4py} process-mining ecosystem, and LLM-to-PDDL
generation. Both surfaces share a governance tier, differing only in
language: neither is ggen-rendered, neither is ledgered in
\texttt{.mfact/artifacts.toml}, and neither carries standing, counts, or
certification data -- their absence from the ledger is, per AGENTS.md,
correct by design, not an omission awaiting correction. \texttt{pylab/}
is explicitly never added to \texttt{just build}, \texttt{just check},
\texttt{just release}, or the artifact ledger; a failure inside it is an
ordinary code-review issue, never \texttt{ARTIFACT\_DRIFT\_REFUSED} or a
release-standing regression, and PYLAB-PRD-ARD-v26.7.7.md is explicit
that nothing in \texttt{pylab/} is claimed to be semantically equivalent
to, or a proof about, any Lean formalization in
\texttt{procint/ProcInt/} -- the same restraint the repository applies to
its own \texttt{ProcInt.Planning.Pddl} module, whose relationship to any
future \texttt{pddl-plus-parser}-based Python experiment is, in the PRD
document's own words, ``two independent things that happen to model
similar concepts,'' not a proven-equivalent pair, absent an actual
correspondence proof this repository does not currently attempt.

\texttt{pylab/}'s tool inventory, researched and wired via \texttt{uv
add} rather than direct \texttt{pyproject.toml} edits, comprises five
research libraries plus the FastMCP server framework used to expose
them. \textbf{TPOT2} (merged into mainline \texttt{tpot} v1.1.0)
provides genetic-programming AutoML, evolving scikit-learn pipelines as
DAGs under multi-objective selection, and is the direct model for the
tactic-search loop of Section~\ref{sec:ai-genetic-search}. \textbf{pm4py}
is the foundation process-mining library and the confirmed reference
implementation of both OCEL~2.0 and Improved Token-Based Replay, serving
as a dependency for the two libraries that follow it. \textbf{powl}
provides POWL~v2 discovery and Petri-net/BPMN/POWL conversion, hard-
depending on \texttt{pm4py}; at the time of the PRD document, a real
version incompatibility existed between \texttt{powl==2.3.7} and
\texttt{pm4py==2.2.32} (a missing internal \texttt{pm4py} module path
\texttt{powl} expects), handled in \texttt{pylab/}'s wrapper by a lazy
import and a structured error rather than a crash. \textbf{ocpa}
supplies object-centric process analysis over OCEL~2.0 data, including
sqlite import and object-centric Petri-net discovery. \textbf{pddl-plus-parser}
parses PDDL+, a continuous/hybrid superset of the PDDL~3.1 fragment
already covered by the repository's Rust \texttt{pddl} crate, and is
retained as a genuinely complementary parser for benchmark sourcing and
cross-validation rather than a duplicate.

Two decisions recorded in the PRD document are worth treating as
governance case studies in their own right, because both concern tools
that were seriously considered and then not adopted, for reasons the
document is careful to record so the decision is not re-litigated.
\textbf{SpiffWorkflow} was considered as a live POWL~v2 execution engine
and explicitly rejected: it is fundamentally BPMN-token-based, compiling
sequence flows and gateway types into a \texttt{TaskSpec} tree, and
feeding it a POWL partial-order block would require a lossy graph
transformation, since arbitrary partial orders do not reduce cleanly to
nested AND/XOR gateways without auxiliary dummy tasks. Since the
sibling \texttt{bcinr-powl} crate already implements POWL~v2 execution
natively with an exact op-kind semantics, there was no reason to adopt a
lossy BPMN intermediary in its place. \textbf{l2p}, an LLM-to-PDDL tool
generating domain and problem PDDL from natural language, was scoped in
the PRD document's first revision as a candidate front end
(\texttt{breed\_l2p}, with a planned FastMCP wrapper); the revision
actually applied instead built a general-purpose FastMCP server exposing
read-only Lean/Lake/\texttt{just} introspection and one tool per
\texttt{pylab} library, superseding the l2p-specific wrapper. AGENTS.md's
governing rule since goes further than that supersession: \texttt{l2p} is
to be deprecated outright -- not installed, referenced, or used as a
planning library -- a stricter position than the PRD document's own
``decided: built with FastMCP'' framing for it. This chapter treats the
later, stricter AGENTS.md rule as authoritative and the PRD document's
earlier framing as superseded history, which is itself a small instance
of the governance discipline this chapter is about: a decision recorded
once is not immutable, and a later, more restrictive rule in AGENTS.md
governs current and future work regardless of what an earlier planning
document scoped.

\section{The Standing-Guard Scanner: Read-Only Gap Detection}
\label{sec:ai-standing-guard}

\texttt{pylab/src/mpops/standing\_guard/server.py} is itself an instance
of the \texttt{pylab/} surface just described: an unledgered,
hand-authored FastMCP server, described in its own module docstring as
``read-only integrity, correctness, and guardrail scan check tools.'' It
exposes a single tool, \texttt{scan}, which runs eight independent gap
checks over the repository and returns each finding as a typed record
carrying a gap class, a severity, a refusal code, the path or target
affected, observed evidence, an expected value, an actual value, and a
recommended action -- itself a small, non-Lean echo of the typed-refusal
discipline of Section~\ref{sec:mfact-candidates}. The eight checks, in
the order the scanner runs them, are:

\begin{enumerate}
\item \textbf{Sorry theorem promotion.} The scanner parses every TTL
  fragment under \texttt{packs/lean-math-pack/fragments/} plus
  \texttt{ontology.ttl} for declarations whose \texttt{procint:status} is
  \texttt{"proven"}, compiles a single scratch Lean file importing each
  declaration's module and running \texttt{\#print axioms} on each
  proven name via \texttt{lake env lean --stdin}, and flags
  \texttt{SORRY\_THEOREM\_PROMOTED} if the kernel output for that name
  mentions \texttt{sorryAx}, or if the name is unknown to the Lean
  environment at all.
\item \textbf{Ledger drift.} The scanner walks every artifact recorded
  in \texttt{.mfact/artifacts.toml}, recomputes a BLAKE3 hash of the file
  on disk (falling back to the \texttt{b3sum} CLI if a native hash
  implementation is unavailable), and flags \texttt{ARTIFACT\_MISSING} if
  the file is absent, \texttt{ARTIFACT\_DRIFT\_REFUSED} if the computed
  hash disagrees with the ledgered hash, and \texttt{ORPHAN\_ARTIFACT\_REFUSED}
  if \texttt{git ls-files} or \texttt{git status --porcelain} reports the
  ledgered file as untracked.
\item \textbf{Orphan artifact scan.} The scanner globs
  \texttt{release/**/*.json} and \texttt{paper/**/*.tex} (excluding
  \texttt{procint/Playground/**}, \texttt{pylab/**}, and
  \texttt{paper/main.tex} itself) for files absent from the ledger, and
  flags \texttt{ORPHAN\_ARTIFACT\_REFUSED} when such a file's content
  looks standing-bearing -- containing any of \emph{proven}, \emph{stated},
  \emph{certified}, \emph{sorry}, \emph{hash}, \emph{audit},
  \emph{quadrature}, or \emph{verif} -- or when it already lives under
  \texttt{release/} or \texttt{paper/} regardless of content.
\item \textbf{regen-check coverage gap.} The scanner extracts the body of
  the justfile's \texttt{regen-check:} recipe and checks, for each
  ledgered artifact, whether its declared producer -- \texttt{ggen}, a
  script path, or another named command -- is textually referenced
  inside that recipe body, flagging \texttt{REGEN\_CHECK\_COVERAGE\_GAP}
  as a warning when it is not: a producer that only runs under some
  other, separately invoked recipe is a coverage gap in
  \texttt{regen-check}, not a convenience.
\item \textbf{Correspondence binding.} The scanner reads
  \texttt{release/verif-receipt.json} and, for each recorded obligation,
  flags \texttt{STALE\_PROOF\_BINDING} if its \texttt{aeneasDecl} field is
  the literal string \texttt{"TBD"}, or if its status is
  \texttt{PROVEN} but the corresponding Lean proof file under
  \texttt{dist/verif/lean/.../Corr/} is missing, or its comment-stripped
  body does not reference the declared \texttt{aeneasModule}. This check
  is developed further, alongside the correspondence rail proper, in
  Chapter~\ref{ch:correspondence}.
\item \textbf{Tag ancestry.} The scanner runs \texttt{git rev-parse} and
  \texttt{git merge-base --is-ancestor} to confirm that the tag
  \texttt{v26.7.7-procint-certified} is an ancestor of \texttt{HEAD} and,
  when \texttt{release-manifest.json} carries a \texttt{runIdentifier},
  that the identifier is itself an ancestor of the tag, flagging
  \texttt{TAG\_ANCESTRY\_FAIL} otherwise -- the mechanized instance of the
  core-identity protection developed in
  Section~\ref{sec:mfact-core-identity}.
\item \textbf{Untracked ontology fragments.} The scanner globs
  \texttt{packs/*/fragments/*.ttl} and flags
  \texttt{UNTRACKED\_ONTOLOGY\_FRAGMENT} for any fragment file git
  reports as untracked, closing the gap between ``committed source
  tree'' and ``what \texttt{cat fragments/*.ttl} actually concatenates.''
\item \textbf{Prose/paper consistency.} The scanner reads
  \texttt{paper/main.tex} line by line and paragraph by paragraph. Two
  rules are self-contained per line: a stale bare count of 145 or 318
  (the pre-\TotalDecls totals reported in STANDING.md's own historical
  ladder) flags \texttt{STALE\_PAPER\_PROSE\_COUNT}, and a phrase
  matching ``Aeneas proves/verified/checked/certified'' flags a
  \texttt{PROSE\_LINT\_VIOLATION} in favor of the more precise ``Aeneas
  extracts'' and ``Lean proves.'' Three further rules are scoped to the
  whole paragraph rather than a single line, to avoid false positives
  from LaTeX's hand-wrapped line breaks splitting a trigger word from its
  qualifier: a bare \emph{chain} or \emph{hash} must be accompanied
  somewhere in the same paragraph by a qualifying phrase such as
  ``receipt chain'' or ``content hash''; a bare \emph{prove}/\emph{proof}/
  \emph{proven} must be accompanied by a formal-context word such as
  \emph{Lean}, \emph{kernel}, or \emph{axiom}; and a totality adverb --
  \emph{entirely}, \emph{completely}, \emph{fully}, \emph{absolutely},
  \emph{always}, \emph{never} -- must be accompanied by an explicit scope
  qualifier. Each unaccompanied occurrence is reported as a
  \texttt{PROSE\_LINT\_VIOLATION} warning at the specific line the
  trigger word occurs on.
\end{enumerate}

The scanner is deliberately read-only: it reports findings, it does not
repair the repository, and its own existence inside \texttt{pylab/} means
it is subject to exactly the same governance tier -- no standing claim,
no ledger entry -- as the surface it audits. It is, in effect, the
Section~\ref{sec:ai-doctrine} doctrine applied reflexively: automated
tooling checking the pipeline's own discipline is itself treated as an
untrusted, unledgered producer of findings, not as a certifying
authority in its own right.

\section{The 2026-07-07 Five-Rail Audit: A Case Study in Self-Correction}
\label{sec:ai-guardrail-audit}

The standing-guard scanner's eight checks did not arise from abstract
risk analysis. A five-rail audit conducted on 2026-07-07, recorded in
\texttt{pylab/docs/jira/26.7.7/tickets/ticket\_013\_v26\_7\_7\_gap\_audit.md},
found that a sorry-backed theorem had been ledgered with
\texttt{procint:status "proven"} in the TTL catalog, with no guard in
place that would have caught it, alongside several related silent-drift
failure modes. AGENTS.md records the resulting structural rules under a
``Guardrails (post-v26.7.7-audit)'' heading, framed explicitly as making
each specific failure mode \emph{structurally impossible}, not merely
discouraged, and each rule maps onto one of the scanner checks of
Section~\ref{sec:ai-standing-guard}:

\begin{itemize}
\item \textbf{No status promotion without a checked guard.} Any TTL
  fragment setting \texttt{procint:status "proven"} on a theorem must
  have a corresponding entry in \texttt{release/gates.json}, or an
  equivalent builder check, that mechanically verifies \texttt{\#print
  axioms} shows no \texttt{sorryAx} for that declaration; a status
  literal in TTL alone is never sufficient, and a human or agent setting
  \texttt{"proven"} without a passing mechanical check is
  \texttt{STATED\_PROMOTED\_TO\_PROVEN}. This rule's mechanized instance
  is the scanner's sorry-theorem-promotion check
  (Section~\ref{sec:ai-standing-guard}, item~1). A named companion
  vocabulary applies specifically to the five deliberately stated
  workflow-countermodel declarations introduced in
  Section~\ref{sec:mfact-boundary}: their status key is
  \texttt{WFNET\_INFINITE\_TRANSITION\_COUNTERMODEL}, the guard that must
  never let them be promoted is named \texttt{countermodel\_not\_promoted},
  and the refusal that guard fires is \texttt{COUNTERMODEL\_PROMOTION\_REFUSED}
  -- language that exists precisely so these five declarations are never
  mistaken for unfinished work awaiting the same promotion path as an
  ordinary stated theorem.
\item \textbf{regen-check must not have untracked-file blind spots.}
  \texttt{git diff --exit-code} only sees tracked files, so a ledgered
  artifact's producer script must itself be invoked by
  \texttt{regen-check}/\texttt{check}, not merely by some separate recipe
  outside that chain; newly created ledgered artifacts must be
  \texttt{git add}ed in the same commit that introduces them, and an
  untracked ledgered file whose disk hash disagrees with
  \texttt{.mfact/artifacts.toml} is \texttt{ORPHAN\_ARTIFACT\_REFUSED} and
  must actually fire, not merely exist as a documented category. This
  rule's mechanized instances are the scanner's ledger-drift and
  regen-check-coverage-gap checks (Section~\ref{sec:ai-standing-guard},
  items~2 and~4).
\item \textbf{Correspondence theorems must reference their extraction.}
  A correspondence obligation reported as \texttt{PROVEN} must have its
  Lean statement actually import and quantify over the extracted or
  generated declaration -- for example an Aeneas \texttt{Generated.*}
  type -- via its declared abstraction function, not merely present a
  same-shaped statement over hand-written types; a receipt field such as
  \texttt{aeneasDecl: "TBD"} is itself a refusal signal that must block
  \texttt{PROVEN} status rather than ship alongside it. This rule's
  mechanized instance is the scanner's correspondence-binding check
  (Section~\ref{sec:ai-standing-guard}, item~5), developed further in
  Chapter~\ref{ch:correspondence}.
\item \textbf{Tag ancestry is part of certification.}
  \texttt{just certify}/\texttt{just release} must verify the release
  commit is an ancestor of, or equal to, the certified git tag before
  reporting \texttt{CERTIFIED\_RELEASE=PASS}; a tag that predates the
  currently rendered artifacts is \texttt{RECEIPT\_RECURSION\_REFUSED},
  not a historical curiosity to note and move past. This rule's
  mechanized instance is the scanner's tag-ancestry check
  (Section~\ref{sec:ai-standing-guard}, item~6), and its target property
  is exactly the core-identity protection of
  Section~\ref{sec:mfact-core-identity}.
\item \textbf{New fragments feeding \texttt{ontology.ttl} must be
  tracked before \texttt{regen-check} runs.} An untracked \texttt{.ttl}
  file silently concatenated into generated output is a source-provenance
  gap distinct from ordinary generated-artifact drift: it means the
  generated output is not reproducible from the committed source tree at
  all. This rule's mechanized instance is the scanner's untracked-fragment
  check (Section~\ref{sec:ai-standing-guard}, item~7).
\end{itemize}

Read together, the audit and its resulting guardrails are the strongest
evidence this thesis can offer, internal to the repository itself, that
the untrusted-producer doctrine of Section~\ref{sec:ai-doctrine} is
enforced rather than merely stated. The failure the audit found was
real, not hypothetical: a status literal reached \texttt{"proven"} in a
committed fragment without the mechanical check that should have gated
it. The response was not a promise to be more careful, but a scanner
whose sorry-theorem-promotion check, run against the repository today,
would have caught exactly that failure, together with six further checks
covering the related silent-drift modes the same audit identified. This
chapter's account of AI-assisted development therefore closes where
Chapter~\ref{ch:mfact} opened: standing is never inferred from the fact
that a candidate -- human-authored, agent-authored, or produced by a
genetic search over tactic sequences -- made contact with the
repository. It is conferred only by an admission step that a later,
independent check can re-run and, when it fails, is designed to catch.

\section{Summary}

This chapter examined the pipeline's own construction under the same
discipline the pipeline enforces on \textsf{procint}'s mathematical
content. Generative-AI tooling is treated, by explicit disclosure and by
mechanism, as an untrusted candidate producer whose output enters the
trusted base only through kernel admission, sorry census, axiom audit,
and release-gate checks -- never through the fact of having been
produced by a more or less capable model. The genetic tactic search tool
instantiates this doctrine for automated proof search specifically:
fitness is kernel accept/reject, never an LLM's self-report, and
promotion from a found candidate to an admitted theorem remains a
separate, human-reviewed step outside the search script's own reach.
\texttt{pylab/} and \texttt{procint/Playground/} provide governed spaces
for hand-authored experimentation -- including the tool inventory and
the SpiffWorkflow and l2p decisions recorded above -- that never acquire
release standing by virtue of existing. And the standing-guard scanner,
together with the five-rail audit and its resulting AGENTS.md guardrails,
is this project's own record of a real gap in that discipline and the
mechanized, read-only check subsequently built to close it. Chapter~\ref{ch:evaluation}
now turns to consolidating the quantitative release evidence this and
earlier chapters have referenced piecemeal into a single evaluation
treatment.
